\documentclass[a4paper,11pt]{article}
%% Compiles both from this folder (pdflatex tp2.tex) and from the root of the project
%% (Overleaf: Menu > Main document = tp2_bagging/tp2.tex).
\makeatletter
\IfFileExists{tp2_bagging/tp2.tex}{\gdef\tpdir{tp2_bagging/}\gdef\tpcommon{common/}}{\gdef\tpdir{}\gdef\tpcommon{../common/}}
\edef\input@path{{\tpdir}{\tpcommon}}
\makeatother
\input{tpsetup}
\renewcommand{\polytitle}{Statistical Supervised Machine Learning -- TP 2}
\graphicspath{{figures/}{tp2_bagging/figures/}{../common/logos/}{common/logos/}}

\begin{document}
\sloppy
\sheethead{TP 2 -- Bagging and random forests}{%
Second of five practical sessions (1: protocol, 2: bagging, 3: boosting, 4 and 5: PAC-Bayes). At the end of TP~1, a
majority vote of trees learned on bootstrap samples was much better than a single tree. In this session we understand
why. We first measure the instability of decision trees and decompose their error into bias and variance; we then
implement bagging, use its out-of-bag error, and check that bagging reduces the variance but not the bias, so that it
helps unstable models only. Finally, random forests decorrelate the trees, and we observe the trade-off between the
accuracy of the voters and their diversity. We end by reproducing the ``puzzling experiment'' that opens Part~IV of the
lecture notes (Section~12.1), with the notation of the PAC-Bayesian analysis of majority votes (TP~4 and~5). The session lasts \textbf{three hours}; Part~\ref{p2:far} contains the difficult questions.}

\begin{sessionbox}[About this session -- TP 2: Bagging and random forests]
\textbf{Goal.} By the end of the session you should be able to explain why trees have a high variance, implement
bagging and compute its out-of-bag error, say which models benefit from bagging, implement a random forest and explain
the role of the number of features drawn at each node.\par
\textbf{To reread before the session.} Section~8.7 (decision trees), Section~10.1 (error decomposition) and Section~10.2
(bagging and random forests) of the lecture notes. For Part~\ref{p2:rf}, have a look at Section~12 (a vote is a distribution
over voters, Definition~12.1) and Section~15.4 (joint error and disagreement, Definition~15.4 and Proposition~15.5).\par
\textbf{Timing (3 hours).}\quad
\begin{tabular}[t]{ll}Part~1 & 30 min\\
Part~2 & 45 min\\
Part~3 & 35 min\\
Part~4 & 60 min\\
Part~5 \hard & 10 min, or at home\end{tabular}\par
\textbf{To hand in.} Your notebook, with the figures of Parts~1, 2 and~4 and the tables of Parts~2 to~4.\par
\textbf{Tools.} \code{numpy}, \code{matplotlib}; the trees of \code{scikit-learn} (\code{DecisionTreeClassifier},
\code{DecisionTreeRegressor}) as base learners. Bagging and random forests are implemented by hand, and compared with
\code{BaggingClassifier} and \code{RandomForestClassifier}.
\end{sessionbox}

\medskip
\noindent\textbf{Data.} \code{wdbc}, \code{pima}, \code{australian} and \code{spambase} ($4597$ e-mails, $57$ features),
loaded with \code{load(name)} of \code{common/tptools.py}, labels in $\{-1,+1\}$, with $70\%/30\%$ stratified
training/test splits (\code{random\_state=seed}); and the regression problem $Y=\sin(2\pi X)+\varepsilon$, $X$ uniform on
$[0,1]$, $\varepsilon\sim\mathcal N(0,0.3^2)$. The protocol of TP~1 applies: the test set is used only to report errors.
The file \code{tp2\_starter.py} contains the skeletons.

\begin{tip}
Store the predictions of all the base models in a $B\times n_{\text{test}}$ matrix \code{V} of $\pm1$ values: the vote of
the first $B$ models is \code{np.sign(np.cumsum(V, 0))} (decide what to do with ties), the mean error of a model is
\code{np.mean(V != y)}, and many quantities of this session are one line of \code{numpy} away.
\end{tip}

%% =====================================================================
\tppart{Trees are unstable (Sections 8.7 and 10.1)}\label{p2:var}

\begin{enumerate}
\item On \code{pima} (split \code{seed=0}), learn two trees on two bootstrap samples of the training set ($m$ indices
      drawn with replacement) and compute the proportion of test points on which they disagree.\\
      Repeat $20$ times, for
      maximal depths $1$, $2$, $5$ and unlimited.\\
      Compare with the test error of a tree.
\item \textbf{Bias and variance.} For the regression problem, draw $R=200$ training samples of size $m=100$, learn a
      regression tree of maximal depth $D$ on each, and estimate on a regular grid of $[0,1]$ the squared bias
      $\overline{(\bar f_D(x)-f(x))^2}$ and the variance $\overline{\mathrm{Var}_r\hat f^{(r)}_D(x)}$, where $\bar f_D$ is the
      average of the $R$ trees and the bar denotes the average over the grid.\\
      Plot them for $D\in\{1,\dots,12\}$.\\
      Recall
      why the expected test MSE is $\text{bias}^2+\text{variance}+\sigma^2$.
\end{enumerate}

%% =====================================================================
\tppart{Bagging from scratch (Section 10.2)}\label{p2:bag}

\begin{enumerate}
\item Show that a bootstrap sample of size $m$ contains on average a proportion $1-(1-\frac1m)^m\approx1-e^{-1}\approx63.2\%$
      of distinct examples.\\
      Check it numerically for $m=537$.
\item Write \code{bagging(Xtr, ytr, Xte, B, rng)} which learns $B$ fully grown trees, each on a bootstrap sample, and
      returns the $B\times n_{\text{test}}$ matrix of their predictions.\\
      On \code{wdbc}, \code{pima} and \code{spambase}
      ($B=200$, $3$ splits), compute the test error of one full tree learned on the whole training set, the mean test error
      of a bagged tree, and the test error of the vote (ties: majority class).\\
      Plot the error of the vote of the first $B$
      trees as a function of $B$.
\item \textbf{Out-of-bag error.} Each training example is \emph{out of bag} for about $37\%$ of the trees.\\
      Predict each
      training example by the vote of these trees only, and compute the out-of-bag (OOB) error.\\
      Compare it with the test
      error.\\
      Why is it a valid estimate of the risk of the vote, without any validation set?
\item Compare with \code{BaggingClassifier(DecisionTreeClassifier(), n\_estimators=200)}.
\item Does the vote overfit when $B$ increases?\\
      How many trees are needed?
\end{enumerate}

%% =====================================================================
\tppart{What does bagging reduce? (Section 10.1)}\label{p2:what}

\begin{enumerate}
\item For the regression problem ($R=100$ training samples of size $100$), compare the squared bias and the variance of a
      single regression tree and of a bagged average of $50$ regression trees, for depths $2$, $4$ and unlimited.
\item On \code{pima} and \code{spambase} ($3$ splits), compare the test error of a model and of the vote of $50$ bagged
      copies of it, for: a fully grown tree, a tree of depth $2$, $15$-NN and a logistic regression (both with
      standardization).\\
      Which models benefit from bagging?\\
      Explain.
\end{enumerate}

%% =====================================================================
\tppart{Random forests: decorrelating the trees (Sections 10.2, 12 and 15.4)}\label{p2:rf}

A random forest is bagging in which, at each node of each tree, the best split is searched among $m_{\text{try}}$
features drawn at random instead of all the $d$ features (\code{DecisionTreeClassifier(max\_features=mtry)});
$m_{\text{try}}=d$ is bagging.

\noindent\textbf{Notation of Part IV.} The vote of $B$ trees with equal weights is the majority vote $\BQ$ of the uniform
distribution $Q$ on the trees (Section~12.2). For an example $(x,y)$, $W_Q(x,y)$ is the proportion of trees that err on
it; the vote errs when $W_Q\geq\frac12$. The mean error of the trees is the \emph{Gibbs risk}
$\RS(\GQ)=\widehat\E[W_Q]$ (Definition~12.1). For two trees $h,h'$ drawn independently from $Q$ (possibly the same
one), the \emph{joint error} $\hat e(Q)$ is the probability that both err and the \emph{disagreement} $\hat d(Q)$ the
probability that they predict differently (Definition~15.4).

\begin{enumerate}
\item For $m_{\text{try}}\in\{1,2,3,5,8\}$ on \code{pima} and $\{1,3,7,15,30,57\}$ on \code{spambase}, learn forests of
      $100$ trees and compute (a) the mean test error of one tree, (b) the mean \emph{disagreement} between two distinct
      trees (the proportion of test points on which they predict differently, averaged over the pairs), (c) the test error
      of the vote.\\
      Plot the three quantities against $m_{\text{try}}$ and comment.
\item Show that, with $v(x)=\frac1B\sum_bh_b(x)$, the mean disagreement over the $B(B-1)$ ordered pairs of distinct trees is
      $\frac{B}{B-1}\,\widehat\E_x\big[\frac{1-v(x)^2}{2}\big]$.\\
      Show that the disagreement $\hat d(Q)$ of Part IV, which averages over all the $B^2$ ordered pairs (including
      $h=h'$), is $\widehat\E_x\big[\frac{1-v(x)^2}{2}\big]=\widehat\E\big[2W_Q(1-W_Q)\big]$, that $\hat e(Q)=\widehat\E[W_Q^2]$, and deduce
      the identity $\RS(\GQ)=\hat e(Q)+\frac12\hat d(Q)$ (Proposition~15.5).
\item On the four datasets, compare a single tree whose depth is tuned by $5$-fold cross-validation (TP~1) with
      \code{RandomForestClassifier(500, max\_features="sqrt", oob\_score=True)}.\\
      Compare the OOB error of the forest with
      its test error.
\item \textbf{The puzzling experiment of Section~12.1.} On the \code{digits} dataset of \code{scikit-learn} (task: is the
      digit smaller than $5$?), split the data with \code{train\_test\_split(X, y, test\_size=0.3, random\_state=0)}
      ($1257$ training and $540$ test digits) and grow $100$ trees, each on a bootstrap sample of \emph{half} the size of the
      training set, with \code{max\_features="sqrt"}.\\
      Plot, as a function of $n$, the test error of the vote of the first $n$ trees, their Gibbs risk and their
      disagreement; plot the histogram of the test errors of the $100$ trees.\\
      Compute $\hat e(Q)$ and $\hat d(Q)$ for the $100$ trees on the test set and check the identity of question~2.\\
      Why is the vote so much better than its trees?
\end{enumerate}

%% =====================================================================
\tppart{Going further \hard}\label{p2:far}

\begin{enumerate}
\item \textbf{Variance of an average.} Let $Z_1,\dots,Z_B$ be identically distributed with variance $\sigma^2$ and
      pairwise correlation $\rho\geq0$.\\
      Show that $\mathrm{Var}\big(\frac1B\sum_bZ_b\big)=\rho\sigma^2+\frac{1-\rho}{B}\sigma^2$.\\
      What does this formula say about bagging when $B\to\infty$, and about random forests?
\item Check it numerically: on \code{pima}, at $40$ fixed test points, compute the average of the predictions ($\pm1$) of
      $B$ bagged full trees, for $60$ training sets (random halves of the training set), and estimate its variance for
      $B\in\{1,5,25,100\}$.\\
      Estimate $\rho$ from $B=1$ and $B=100$ and compare the formula with the other values.
\item If $B$ trees erred \emph{independently} with probability $r<\frac12$ on every point, what would happen to the error
      of the vote when $B\to\infty$ (Condorcet's jury theorem, Proposition~15.3)?\\
      Why is this not what we observe?
\end{enumerate}

\end{document}
